\chapter*{Introducción}
\addcontentsline{toc}{chapter}{Introducción}
\pagenumbering{arabic}


\begin{quote}
	While the move from dimension 2 to dimension 3 appears to be the obvious step there is a sense in which one should move from 2 to 4. This comes from the consideration of complex algebraic geometry. For complex dimension 1 this theory was started by Abel and continued by Riemann. For algebraic varieties of complex dimension n the real dimension is 2n. The key figures in the topology of higher-dimensional algebraic varieties were Lefschetz, Hodgex, Cartan and Serre. While general algebraic geometry was one of the major developments of the second half of the 20th century, the topology of real 4-manifolds had a great surprise in store when Simon Donaldson made spectacular discoveries opening up an entirely new area.
	
	\hfill --- James Carlson, \textit{The Poincaré Conjecture}, AMS (2014)
\end{quote}


En la intersección entre el análisis complejo y la geometría diferencial surge una de las ramas más hermosas de la geometría; la \textit{geometría compleja}. Esta rama se encarga de estudiar objetos como variedades complejas, variedades algebraicas, funciones en varias variables, haces vectoriales holomorfos, entre otros. En particular, se hace un estudio profundo de objetos que generalizan las superficies de Riemann en dimensiones mayores; las variedades complejas. Este enfoque constituye un pilar para numerosas teorías contemporáneas, incluyendo la cohomología de Dolbeault y de gavillas, la teoría de Hodge, el estudio de foliaciones holomorfas, la teoría de singularidades y la teoría de cuerdas.

Una variedad compleja es un espacio que localmente se comporta como un abierto del espacio complejo $\cc^n$, pero cuyo comportamiento global puede ser mucho más profundo. Además, se encuentra equipado con funciones de transición holomorfas, lo cual garantiza una compatibilidad entre la geometría local compleja y todo el espacio. Sin embargo, la condición de holomorfia es más rígida que la suavidad real. Por ejemplo, las variedades complejas compactas no admiten funciones holomorfas globales no constantes, y el espacio de secciones globales de un haz vectorial holomorfo siempre tiene dimensión finita, en el caso compacto. La inexistencia de un teorema del encaje de Whitney general, sumado a fenómenos como el teorema de Hartogs, el comportamiento local de las funciones holomorfas y el concepto de funciones meromorfas, nos muestran una teoría esencialmente distinta a la geometría diferencial real. Por esta razón, buscamos herramientas más manejables para abordar las variedades complejas. Una diferencia importante respecto a las variedades suaves reales, y la cual es particularmente relevante para este trabajo, es la cercanía de las variedades complejas con las variedades algebraicas. En particular, estas conexiones con la geometría algebraica son recurrentes y una parte profunda de la teoría.

La \textbf{algebreización de variedades complejas} consiste, de manera informal, en asociar a una variedad compleja una variedad algebraica que la reproduzca como objeto analítico. Esto permite reemplazar un objeto analítico por uno algebraico, ganando acceso a herramientas mucho más rígidas, estructuradas y computacionales. Por ejemplo, en geometría algebraica, las funciones meromorfas del espacio proyectivo son cocientes de polinomios homogéneos del mismo grado. De este modo, podemos reemplazar las complejidades que conlleva la descripción holomorfa (continuidad, diferenciabilidad y convergencia) por problemas puramente algebraicos (ecuaciones polinomiales, estructuras algebraicas, invariantes algebraicos). De esta manera, \textit{algebrizar significa mostrar que la complejidad analítica está gobernada por su estructura algebraica subyacente}.

El objetivo del trabajo es presentar las propiedades básicas de las variedades complejas, incluyendo un estudio breve de las aplicaciones y funciones holomorfas, el concepto de gavilla sobre una variedad compleja, así como presentar ejemplos importantes de estos objetos. Posteriormente, se abordará el problema de encajar una variedad compleja a un espacio proyectivo, para lo cual se introduce el concepto de haz lineal y la idea de amplitud. En el último capítulo, entenderemos y demostraremos el teorema de Chow, apoyándonos en los resultados de Remmert. Este teorema será el eje central de todo el trabajo, pues conecta todos los conceptos estudiados. Finalmente, en las dos últimas secciones vamos a presentar brevemente los conceptos de variedad de Moishezon y Kähler, con el propósito de entender su importancia en el tema de la algebrización. 

A lo largo del trabajo se asumirá familiaridad con los fundamentos de topología general y análisis complejo en una variable. Sin embargo, se incluye un apéndice con resultados sobre funciones holomorfas de varias variables que serán citados a lo largo de todo el trabajo. Asimismo, se hará referencia durante todo el trabajo a conceptos y resultados clásicos del álgebra conmutativa, geometría algebraica y superficies de Riemann, remitiendo al lector a la bibliografía clásica y citando únicamente aquellos resultados técnicos que sean necesarios. Para las dos últimas secciones se utilizarán nociones más avanzadas de la geometría algebraica y geometría diferencial. Así, dichos temas solo serán presentados de manera breve y no se hará un desarrollo exhaustivo.

